\chapter{Metatheory}
\label{ch:metatheory}

The preceding three chapters develop logics; this closing chapter
asks questions \emph{about} them. Soundness, decidability, and the
representability of syntax as an ordinary object are not full
formal developments here — that would be its own book — but each
gets a concrete, checkable instance rather than a purely expository
treatment.

\section{Soundness for the Propositional Fragment}

A textbook soundness theorem usually relates two separately defined
things: a syntactic proof system (axioms and inference rules) and a
semantics (valuations, truth tables), showing that everything
provable in the first is true in the second. This book has not built
a separate proof system for \texttt{K3} — every claim in
Chapter~\ref{ch:nonclassical} was proved directly in Lean's own type
theory, with Lean's kernel as the arbiter. That changes what
soundness needs to mean here, not whether it matters: the relevant
question is not ``does provability track truth,'' but ``does the
Lean encoding of Kleene's connectives actually compute what
Kleene's own truth tables say it should.'' That is a real,
checkable question, and getting it wrong would silently invalidate
everything built on top of it.

\begin{experiment}{Base cases: the encoding matches its own semantics}
\begin{lean}
theorem knot_involutive : ∀ p : K3, p.knot.knot = p := by decide

theorem kor_comm : ∀ p q : K3, K3.kor p q = K3.kor q p := by decide
\end{lean}
Both are properties Kleene's informal truth tables guarantee — 
negation should undo itself, disjunction should not care about
argument order — checked here not by re-reading the definitions but
by exhaustive computation over every case.
\end{experiment}

The sharper point is what \emph{not} to assume. A reader trained on
classical logic might expect an ``unknown'' input to always leave
\texttt{kor}'s result unknown — that intuition is exactly what
Kleene's own semantics denies, and only Kleene's semantics, checked
directly, can settle it:
\begin{lean}
example : K3.kor K3.unknown K3.tt = K3.tt := by decide
\end{lean}
An unknown left disjunct does not prevent a definite answer, because
the right disjunct alone already settles the question — the same
asymmetric rule stated informally back in
Chapter~\ref{ch:nonclassical}, confirmed here rather than merely
restated. Soundness, in a logic with more than two truth values,
means checking claims like this one against the values that logic
actually has, not against classical habits carried over by default.

\section{Decidability}

\begin{experiment}{Propositional validity over \texttt{K3} is decidable}
\begin{lean}
theorem kor_idempotent : ∀ p : K3, K3.kor p p = p := by decide
\end{lean}
\texttt{K3} has exactly three values, so any formula built from
finitely many \texttt{K3}-valued variables using \texttt{knot} and
\texttt{kor} ranges over a finite space of cases — the same
principle that made every \texttt{decide} call in
\textit{Proof and Countermodel} work for \texttt{Bool}, now applied
to a strictly larger but still finite domain. Every genuinely
propositional claim about \texttt{K3} made in this book —
\texttt{excluded\_middle\_fails}, \texttt{knot\_involutive},
\texttt{kor\_comm}, and this one — was already an instance of this
same decidability fact, used silently each time.
\end{experiment}

\section{Syntax as an Object: A Small Gödel-Numbering Demonstration}

The smallest nontrivial formula language — one atom, and negation as
the only connective — already makes the point without needing
anything as involved as Cantor pairing: a formula is either the
atom or a negation of a smaller formula, so its natural number code
can simply be how many negations it carries.

\begin{experiment}{Formulas as numbers, and back again}
\begin{lean}
inductive Formula where
  | atom
  | neg (f : Formula)

def encode : Formula → Nat
  | .atom => 0
  | .neg f => encode f + 1

def decode : Nat → Formula
  | 0 => .atom
  | n + 1 => .neg (decode n)

theorem encode_decode : ∀ n : Nat, encode (decode n) = n := by
  intro n
  induction n with
  | zero => rfl
  | succ k ih =>
    show encode (decode k) + 1 = k + 1
    rw [ih]
\end{lean}
Once \texttt{encode} and \texttt{decode} exist, a formula is no
longer only something the system manipulates — it is data the
system can compute with, store, and inspect, the same way any other
number can be. \texttt{isWellFormed} below is deliberately trivial
for this minimal language, and says so honestly rather than
performing complexity it does not have:
\begin{lean}
def isWellFormed (n : Nat) : Bool := true

theorem isWellFormed_always : ∀ n : Nat, isWellFormed n = true :=
  fun _ => rfl
\end{lean}
Every natural number decodes to some formula here, because
\texttt{atom}/\texttt{neg} is rich enough to demonstrate the
encoding but not rich enough to produce any malformed codes. A
richer language — a second connective taking two subformulas, coded
by an actual pairing function — would make \texttt{isWellFormed}
genuinely selective, rejecting numbers that do not decode to a
matched pair of valid subformula codes. The encoding technique is
the same either way; only the amount of structure being encoded
changes.
\end{experiment}

\section{Where This Leaves the Admissibility Program}

\texttt{Admissible} from Chapter~\ref{ch:admissibility} is decidable
in any concrete instance the same way \texttt{K3} validity is —
given an explicit history, an explicit constraint, and an explicit
fact, checking membership is a finite computation. That is not the
open question. What this book has not done, anywhere, is separate
\emph{syntax} from \emph{semantics} for admissibility the way a
soundness theorem needs: every claim about \texttt{Admissible} was
proved directly against its Lean definition, never against an
independently stated set of inference rules that could in principle
diverge from it. Chapter~\ref{ch:admissibility}'s three structural
facts — the assumption rule, persistence under extension, and the
explicit failure of constraint-monotonicity — are exactly the
shape of axioms a standalone proof system for $\vdash_{\!A}$ would
need to start from. Writing that system down, independently of the
Lean definition it is meant to describe, and then proving it sound
and complete with respect to \texttt{Admissible} the way modal
logics are proved sound and complete with respect to Kripke
semantics, is the natural next step this book deliberately does not
take. It would be a second book's worth of work, not a section's,
and pretending otherwise here would cost more in false confidence
than a frank ``not yet'' costs in incompleteness.

\section{Exercises}

\begin{exercise}
Prove the other half of the encoding bijection,
$\forall f, \mathtt{decode}\ (\mathtt{encode}\ f) = f$, by induction
on \texttt{Formula}, and explain why both halves together are needed
to call \texttt{encode}/\texttt{decode} a genuine bijection rather
than a one-sided embedding.
\end{exercise}

\begin{exercise}
Extend \texttt{Formula} with a binary connective,
\texttt{conj (f g : Formula)}, taking two subformulas rather than
one. Design a pairing-function-based \texttt{encode} for the
extended language, then define \texttt{isWellFormed} so that it
genuinely rejects some natural numbers — unlike the atom/neg-only
version above, where every number was well-formed by construction.
\end{exercise}
